%==============================================================================
%  CAPÍTULO III
%==============================================================================
\chapter{Institucionalidad concursal y auxiliares de la justicia}
\label{cap:institucionalidad}

\fichaCapitulo{Órganos del concurso y su fiscalización.}{Estructura de la \superir{},
categorías A y B de veedores y liquidadores, y control de cuentas.}

%==============================================================================
\bloqueTeorico
%==============================================================================

\subsection{El diseño orgánico del sistema de insolvencia chileno}

El sistema concursal chileno reparte la administración del concurso entre tres tipos de
actores, y comprender esa distribución es la clave para litigar con eficacia. El
\destacar{tribunal} conserva la jurisdicción: dicta las resoluciones, resuelve las
controversias y controla la legalidad del procedimiento. El \destacar{auxiliar} ---veedor en
la reorganización, liquidador en la liquidación--- administra: informa, incauta, realiza,
reparte. Y la \destacar{\Superir{}} fiscaliza: registra a los auxiliares, los supervigila y
los sanciona.

Esta tríada ---juez que juzga, auxiliar que administra, superintendencia que fiscaliza---
sustituyó al esquema anterior, en que el síndico concentraba administración y una cuota de
autoridad cuasijurisdiccional bajo un control más difuso. La separación de funciones no es
un tecnicismo: define a quién dirigirse para cada cosa. Una objeción de legalidad va al
tribunal; una queja por la conducta del liquidador, a la Superintendencia; una discrepancia
sobre la realización, a la junta.

\subsection{Transparencia del mercado del crédito e imparcialidad del auxiliar}

El liquidador administra un patrimonio ajeno en beneficio de una masa que no lo eligió y con
la que no tiene vínculo. Esa posición ---administrar bienes de otros para pagar a terceros---
genera un riesgo estructural de captura: el auxiliar puede ser tentado a favorecer al deudor
que lo propuso, o al acreedor mayoritario que controla su continuidad, en desmedro de la masa.

Todo el estatuto de los auxiliares ---inhabilidades, prohibición de concertación previa,
régimen de honorarios tabulado, fiscalización de la Superintendencia, control de la cuenta
final--- se explica como respuesta a ese riesgo. La imparcialidad del auxiliar no se presume:
se construye mediante controles, y la transparencia del mercado del crédito depende de que
esos controles funcionen. El \art{36} lo expresa al definir la posición del liquidador:
\destacar{representa los intereses generales de los acreedores y los derechos del deudor en
cuanto puedan interesar a la masa}, no los de quien lo nominó.

%==============================================================================
\bloqueNormativo
%==============================================================================

\subsection{Naturaleza jurídica de la SUPERIR}

El \art{331} \parencite{ley-20720} crea la \Superir{} como \destacar{servicio público
descentralizado, con personalidad jurídica y patrimonio propios}, institución autónoma que se
relaciona con el Presidente de la República a través del Ministerio de Economía, Fomento y
Turismo, con domicilio en Santiago sin perjuicio de sus direcciones regionales.

\subsection{Funciones y atribuciones}

El \art{332} le encomienda \destacar{supervigilar y fiscalizar} las actuaciones de los
veedores, liquidadores, martilleros concursales, administradores de la continuación de las
actividades económicas y asesores económicos de insolvencia, además de las funciones que la
ley le entrega en el Capítulo V ---la renegociación de la persona deudora, que la
Superintendencia tramita directamente---.

El \art{337} detalla sus atribuciones y deberes, comenzando por la fiscalización de todos los
entes sujetos a su control.

\subsection{La potestad sancionatoria: el artículo 338}

El \art{338} configura el poder disciplinario del órgano. Los entes fiscalizados que
infrinjan las normas concursales o incumplan las instrucciones de la Superintendencia pueden
ser objeto de:

\begin{enumerate}
  \item Censura por escrito.
  \item \textbf{Multa} a beneficio fiscal de \destacar{1 hasta 1.000 unidades tributarias
  mensuales}.
  \item Suspensión hasta por seis meses para asumir nuevos procedimientos.
  \item \textbf{Exclusión de la nómina} respectiva.
\end{enumerate}

\begin{foco}[La escala sancionatoria condiciona la conducta del auxiliar]
El rango de la multa ---hasta mil UTM--- y la posibilidad de exclusión de la nómina explican
por qué el liquidador diligente prioriza el cumplimiento de las instrucciones de la
Superintendencia (capítulo \ref{cap:normativa-superir}) sobre cualquier otra consideración.
Para el acreedor que enfrenta a un liquidador negligente, la vía disciplinaria ante la
Superintendencia es, a menudo, más eficaz que la judicial.
\end{foco}

\subsection{La Nómina de Veedores y la garantía de fiel desempeño}

El \art{12} detalla las menciones de la \destacar{Nómina de Veedores} que lleva la \superir{}:
nombre, profesión, domicilio y regiones de ejercicio; las \destacar{calificaciones obtenidas en
el examen} de los últimos cinco años (\art{14}); y el número de procedimientos en que ha
intervenido. La nómina es, así, una ficha pública de desempeño consultable antes de proponer o
aceptar a un veedor.

El \art{16} exige a todo veedor mantener en la Superintendencia una \destacar{garantía de fiel
desempeño de \uf{2.000}}, con vigencia mínima de tres años renovable, mientras subsista su
responsabilidad. Sin ella, el veedor no puede asumir nuevos procedimientos. La garantía puede
consistir en boleta bancaria, póliza de seguro u otra que la Superintendencia acepte
(desarrollada en la \ncg{10}, capítulo \ref{cap:normativa-superir}).

\subsection{Requisitos, prohibiciones e inhabilidades de los veedores}

El \art{17} enumera las prohibiciones para ser veedor. No pueden serlo, entre otros, quienes
hayan sido condenados por crimen o simple delito, ni los funcionarios de la Administración del
Estado ni quienes presten servicios ---remunerados o no--- a la Superintendencia, salvo la
excepción de las labores docentes en educación superior.

\subsection{Exclusión de la nómina y su reclamo}

El \art{18} tasa las \destacar{causales de exclusión} de la nómina: haber sido nombrado en
contravención al Título, dejar de cumplir los requisitos del \art{13}, o ---la de mayor peso
disciplinario--- \destacar{adquirir para sí o para terceros bienes de la masa} que administra,
directamente o por interpósita persona relacionada.

Excluido de la nómina, el \art{19} concede al veedor un \destacar{reclamo} ante el juzgado civil
de su domicilio, dentro de diez días desde la notificación por carta certificada, tramitado
conforme al procedimiento sumario del \art{341}.

\begin{foco}[La prohibición de autocontratación es la piedra angular del estatuto ético]
Que el veedor o liquidador no pueda adquirir los bienes que administra ---ni él, ni un tercero
relacionado--- es la regla que previene el conflicto de interés más grave: rematar barato para
comprar barato. Su infracción no es una falta menor: acarrea la exclusión de la nómina. Para el
acreedor que detecta que un allegado del auxiliar se adjudicó un bien, esta es la vía.
\end{foco}

\subsection{Nominación del veedor}

El \art{22} regula el mecanismo. Recibidos los antecedentes del \art{55}, la \superir{} notifica
a los \destacar{tres mayores acreedores} del deudor dentro del día siguiente y por el medio más
expedito. Dentro del \destacar{segundo día}, cada acreedor propone por escrito o correo
electrónico un veedor titular y uno suplente de la nómina vigente. El sistema busca así que sean
los principales acreedores ---no el deudor--- quienes incidan en la designación del auxiliar que
fiscalizará la reorganización.

\subsection{Cese en el cargo del veedor}

El \art{23} dispone que el veedor cesa por el término del procedimiento o por cese anticipado,
pero \destacar{su responsabilidad subsiste hasta la aprobación de su cuenta final de
administración}. La regla es paralela a la del liquidador (\art{49} y siguientes) y explica por
qué la cuenta final es la última ventana de control sobre el auxiliar.

\subsection{Cese anticipado en el cargo}

El \art{24} tasa las causales de cese anticipado del veedor: revocación por la junta, remoción
decretada por el tribunal, renuncia aceptada por causa grave, haber dejado de integrar la
nómina, e inhabilidad sobreviniente.

\begin{practica}[La inhabilidad sobreviniente hay que vigilarla durante todo el procedimiento]
El cese por inhabilidad sobreviniente (art. 24 \Nro 5) no es un control de una sola vez: una
relación con un acreedor que emerge a mitad del procedimiento puede inhabilitar al auxiliar y
obligar a la designación del suplente. Conviene revisar las relaciones del veedor no solo al
nominarlo, sino cada vez que ingresa un nuevo acreedor relevante.
\end{practica}

\subsection{Deberes del liquidador: el artículo 36}

El \art{36} define la posición y los deberes del liquidador. Representa judicial y
extrajudicialmente los intereses generales de los acreedores y los derechos del deudor en
cuanto interesen a la masa, y debe especialmente incautar e inventariar los bienes,
liquidarlos y efectuar los repartos conforme a la ley.

\subsection{Honorarios del liquidador y régimen de descuento}

El \art{39} sujeta los honorarios del liquidador a la tabla progresiva por tramos del \art{40},
con naturaleza de remuneración única y de gasto de administración del procedimiento. La
\destacar{Nómina de Liquidadores} publica, conforme al \art{33}, el régimen de descuento de
honorarios que cada liquidador ofrece respecto de esa tabla, junto con estadísticas de su
desempeño ---número de procedimientos, principales acreedores, porcentaje de cuentas finales
aprobadas---.

\begin{foco}[La nómina es una herramienta de \emph{due diligence} del auxiliar]
La información que publica la Nómina de Liquidadores ---descuentos ofrecidos, historial,
tasa de aprobación de cuentas--- permite evaluar al auxiliar antes de proponerlo o de
aceptar su nominación. Un liquidador con baja tasa de cuentas aprobadas es una señal que el
acreedor diligente lee antes de la junta, no después.
\end{foco}

\subsection{La cuenta final de administración: arts.\ 49 a 53}

La cuenta final es el acto de cierre del auxiliar y el mecanismo por el cual la masa controla su
gestión. Su régimen, aplicable al liquidador y ---por remisión--- al veedor, se articula en
cinco preceptos.

\subsubsection{Contenido y oportunidad}

El \art{49} entrega a la \superir{} la fijación de la forma y los contenidos obligatorios de la
cuenta mediante norma de carácter general (la \ncg{7} y la \ncg{27}, capítulo
\ref{cap:normativa-superir}), con observancia de la normativa contable, tributaria y financiera.

El \art{50} fija la \destacar{oportunidad}: el liquidador debe acompañar la cuenta al tribunal y
a la \superir{} \destacar{dentro de treinta días} desde que ocurra una de tres circunstancias:
el vencimiento de los plazos legales de realización, el agotamiento de los fondos o el pago
íntegro de los créditos reconocidos, o el cese anticipado del cargo.

\subsubsection{Rendición}

El \art{51} obliga al liquidador, acompañada la cuenta, a \destacar{citar a junta de acreedores}
para rendirla, explicar su contenido y acreditar la retención del porcentaje de honorarios del
\art{39} \Nro 7. La Superintendencia puede concurrir con derecho a voz.

\subsubsection{Objeción}

El \art{52} legitima al \destacar{deudor, a cualquier acreedor y a la \superir{}} para objetar
la cuenta, dentro de los \destacar{cinco días} siguientes a la junta en que se rindió o debió
rendirse. No deducidas objeciones oportunas, la cuenta se tiene por aprobada.

\subsubsection{Ejecución del rechazo}

El \art{53} regula la ejecución de la resolución que rechaza la cuenta: si ordena al liquidador
\destacar{restituir dinero a la masa}, dispone de treinta días ---prorrogables--- desde que la
resolución quede firme para cumplir, bajo apercibimiento de las medidas de ejecución que la
norma detalla.

\begin{foco}[Los cinco días para objetar la cuenta son fatales]
La aprobación de la cuenta cierra la responsabilidad del auxiliar (arts. 23 y 49). El acreedor o
el deudor que sospechan una gestión negligente ---bienes malvendidos, gastos injustificados---
tienen \emph{cinco días} desde la junta de rendición para objetar. Dejar pasar ese plazo
consolida la cuenta y extingue la vía de la responsabilidad. Es, junto con la comparecencia a la
junta constitutiva, el momento de control que no puede perderse.
\end{foco}

%==============================================================================
\bloquePractico
%==============================================================================

\subsection{Guía de obtención de la nominación del veedor o liquidador}

\begin{enumerate}
  \item \textbf{Solicitud de nominación.} Presentación de la solicitud de inicio del
  procedimiento voluntario o forzoso ante el tribunal competente.

  \item \textbf{Oficio a la \superir{}.} El tribunal despacha oficio por vía
  interconectada para que la Superintendencia emita el Certificado de Nominación.

  \item \textbf{Mecanismo de selección.}
  \begin{itemize}
    \item \emph{Sorteo público:} procedimiento por defecto, regulado en la \ncg{17} \citeyearpar{ncg-17},
    aplicable a los procedimientos simplificados y ante la rebeldía del deudor en
    proponer terna.
    \item \emph{Propuesta de acreedores:} acuerdo de las tres mayores mayorías de
    acreedores concurrentes, en los procedimientos ordinarios.
  \end{itemize}

  \item \textbf{Aceptación del cargo.} El veedor o liquidador debe comparecer a aceptar
  dentro del plazo legal de tres días hábiles, bajo apercibimiento de exclusión del
  registro, rindiendo la garantía de fiel desempeño.
\end{enumerate}

\verificar{Confirmar el plazo de aceptación y el monto/forma de la garantía de fiel
desempeño en el texto vigente.}

\begin{practica}[Control de inhabilidades antes de la nominación]
Conviene revisar de antemano las relaciones del veedor o liquidador propuesto con los
principales acreedores. Una inhabilidad sobreviniente detectada en junta obliga a repetir
la nominación y consume semanas de tramitación.
\end{practica}

%==============================================================================
\bloqueJurisprudencial
%==============================================================================

\subsection{Límites de la potestad sancionatoria de la \superir{}}
\pendiente{Sistematizar jurisprudencia de la \csup{} y dictámenes de la Contraloría
General de la República.}

\subsection{Responsabilidad civil del liquidador y la cuenta final de administración}

La responsabilidad civil del liquidador por la devaluación negligente de activos o la
administración descuidada del patrimonio se persigue tras la rendición de la \destacar{cuenta
final de administración}. Conforme a los \art{23} y \art{49}, esa responsabilidad subsiste
hasta la aprobación formal de la cuenta por el tribunal, y hasta entonces los acreedores pueden
entablar demandas incidentales de cobro de perjuicios o exigir restituciones de fondos.

El \art{49} fija que el liquidador dispone de un plazo de \destacar{treinta días} para
presentar la cuenta final desde el vencimiento de los plazos de realización o el agotamiento de
los fondos. Omitida esa presentación, el \art{35} habilita la persecución directa de su
responsabilidad civil, sin perjuicio de las sanciones disciplinarias de suspensión e
inhabilitación que determine la \superir{} en ejercicio de sus facultades fiscalizadoras
(\art{338}).\verificar{Individualizar los fallos que han determinado responsabilidad civil
personal del liquidador por bienes no realizados oportunamente o dañados bajo su custodia.}

\begin{foco}[La cuenta final es la última ventana de control]
Mientras la cuenta final no esté aprobada, la responsabilidad del liquidador sigue viva. Para
el acreedor que sospecha una realización negligente ---un bien rematado a precio vil, una
demora que devaluó el activo---, la objeción de la cuenta final (\ncg{18}) y la demanda
incidental de perjuicios son la vía; dejar aprobar la cuenta sin objetarla cierra esa puerta.
\end{foco}

\subsection{Los honorarios del liquidador: la tabla del artículo 40}

El \art{40}, modificado por la \Lreforma{}, sujeta los honorarios del liquidador a una
\destacar{tabla progresiva por tramos} calculada sobre los fondos efectivamente distribuidos en
cada reparto, con porcentaje decreciente a medida que aumenta el recupero y una remuneración
mínima garantizada.

\begin{table}[htbp]
\centering
\small
\caption{Honorarios del liquidador (tabla progresiva del artículo 40)}
\label{tab:honorarios-liquidador}
\begin{tabularx}{\textwidth}{@{}>{\raggedright\arraybackslash}p{0.44\textwidth} p{0.20\textwidth} X@{}}
\toprule
\textbf{Tramo de distribución acumulada} & \textbf{Porcentaje} & \textbf{Observaciones} \\
\midrule
Hasta \uf{2.000} & 20{,}00\,\% & Remuneración mínima garantizada: \uf{30} \\
\addlinespace
Sobre \uf{2.000} y hasta \uf{10.000} & Escala progresiva intermedia decreciente & — \\
\addlinespace
Sobre \uf{260.000} y hasta \uf{520.000} & 1{,}75\,\% & — \\
\addlinespace
Sobre \uf{520.000} y hasta \uf{1.000.000} & 1{,}50\,\% & — \\
\addlinespace
Sobre \uf{1.000.000} & 1{,}00\,\% & — \\
\bottomrule
\end{tabularx}
\end{table}

\verificar{Completar los tramos intermedios de la tabla del artículo 40 con sus porcentajes
exactos, y confirmar el monto de la remuneración mínima garantizada y quién la solventa.}

\begin{foco}[Por qué el liquidador prioriza el recupero temprano]
La progresión decreciente ---20\,\% en el primer tramo, 1\,\% sobre el millón de UF--- concentra
el incentivo del liquidador en los primeros repartos. Comprender la estructura de honorarios
permite al acreedor anticipar el comportamiento del auxiliar y, al deudor, dimensionar el costo
de administración que soportará la masa.
\end{foco}
